\documentclass[11pt]{article}
\usepackage{mathblatt}
% Probe 08.10.2026: „Die lange Seite“ in drei Durchgängen, Durchgang 2 (Wiederholung), von Hand gesetzt
% nach aufgabenbank bau/bauregeln.md (Stand 08.10.2026) und mathe-nachhilfe plan.md § 5 „Richtung Durchgänge“.
% Wiederholung: ohne Vorstufen, Hilfe nur am Anfang (Nr. 2), kurzes Päckchen, krumme Zahlen; gleiche Sorten, andere Varianten.
% Gemeinsame Präambel der drei Durchgänge, übernommen aus paket-pythagoras-2/2-blatt.tex (Probe 08.10.2026).
\makeatletter
\setlength{\mbpfbe}{0mm}% keine Punkte (Bauregel 6.4)
% eingerückt (Bauregel 4.7, 6.6): \gEin Einzug, \gSchrift eine Stufe kleiner
\newlength{\gEin}\setlength{\gEin}{0pt}
\let\gSchrift\relax
\renewcommand{\pfkreuzzeile}[1]{\par\noindent\hangindent3.2em\hangafter1\hspace*{1.6em}$\square$\ #1\par}
\newcommand{\gkreuz}[1]{$\square$\ #1\hspace{2em}}
% Bauregel 6.7: links, was man ansieht, rechts, was man tut; Spaltengrenze überall an derselben Stelle
\newsavebox{\gbox}\newlength{\gLinks}
\newcommand{\gzwei}[2]{\par\smallskip\noindent\sbox{\gbox}{#1}%
  \setlength{\gLinks}{\dimexpr0.5\textwidth-\gEin-\mbpfnr\relax}%
  \ifdim\wd\gbox>\dimexpr\gLinks-4mm\relax
    \usebox{\gbox}\par\smallskip\noindent\begin{minipage}{\linewidth}#2\end{minipage}%
  \else
    \begin{minipage}[t]{\gLinks}\vspace{0pt}\usebox{\gbox}\end{minipage}%
    \begin{minipage}[t]{\dimexpr\linewidth-\gLinks\relax}\vspace{0pt}\raggedright #2\end{minipage}%
  \fi\par}
\RenewDocumentEnvironment{pfaufg}{m m m}{%
  \begin{lrbox}{\mbaufgabebox}\begin{minipage}[t]{\linewidth}%
  \gSchrift\noindent\hspace*{\gEin}\mbpfrandmarke{#1}%
  \begin{minipage}[t]{\mbpfnr}\strut\textbf{#2}\end{minipage}%
  \begin{minipage}[t]{\dimexpr\linewidth-\gEin-\mbpfnr\relax}\raggedright\strut\ignorespaces}%
 {\end{minipage}%
  \end{minipage}\end{lrbox}%
  \setlength{\mbaufgabehoehe}{\dimexpr\ht\mbaufgabebox+\dp\mbaufgabebox\relax}%
  \par\addvspace{7pt}\Needspace*{\mbaufgabehoehe}%
  \noindent\usebox{\mbaufgabebox}\par\addvspace{7pt}}
\makeatother
% Rechenraum als Karo (Bauregel 6.9): n Rechenschritte = n mal 10 mm
\newcommand{\karo}[1]{\par\smallskip\noindent\begin{tikzpicture}
  \clip (0,0) rectangle ({\linewidth-1mm},{#1*10mm});
  \draw[black!16,line width=0.3pt,step=5mm] (0,0) grid ({\linewidth},{#1*10mm});
  \end{tikzpicture}\par}
\newcommand{\kk}[1]{$\square$\,#1\hspace{0.9em}}
\newcommand{\linie}[1]{\,{\color{black!45}\rule[-1pt]{#1}{0.4pt}}\,}
% Zeile am Ende jedes Durchgangs: Originale klein grau, darüber; dann Vorgänger/Nachfolger (Bauregel 3.4, 6.5)
\newcommand{\blattende}[2]{\par\vfill\noindent\begin{minipage}{\linewidth}{\scriptsize\color{mbgrau}Originale: #1\par}\smallskip
{\footnotesize #2\par}\end{minipage}}
\newcommand{\pkt}{\textperiodcentered{}}

\newcommand{\kennung}{H7Q}
\newcommand{\pz}[3]{\noindent\makebox[41mm][l]{#1}$c #2$\linie{14mm}#3\par\smallskip}

\begin{document}
\pfheftstil
\fancyfoot[C]{\parbox[t]{\textwidth}{\mbfhbox\par\vspace{1.5mm}{\footnotesize\color{mbgrau}{\scriptsize \kennung}\hfill\thepage}}}
\pfheftkopf{Die lange Seite \pkt{} Hypotenuse}{P10 \pkt{} 2}

% 1 Satz in Worten (Kette 2 s7 v2)
\begin{pfaufg}{}{1.}{}
Kreuze die Aussage an, die in jedem rechtwinkligen Dreieck gilt.
\pfkreuzzeile{Die Längen der Katheten ergeben zusammen die Länge der Hypotenuse.}
\pfkreuzzeile{Das Hypotenusenquadrat ist doppelt so groß wie ein Kathetenquadrat.}
\pfkreuzzeile{Die Flächen der zwei Kathetenquadrate ergeben zusammen die Fläche des Hypotenusenquadrats.}
\fusshilfe{\mbox{1:~die dritte Aussage}}
\end{pfaufg}

% 2 Gleichung aufstellen (Kette 2 s6 v1): Hypotenuse liegt unten; die einzige Hilfe des Durchgangs steht am Anfang
\begin{pfaufg}{}{2.}{}
Schreibe für dieses Dreieck die Gleichung nach dem Satz des Pythagoras auf.
Suche zuerst die Hypotenuse.
\gzwei{\begin{tikzpicture}[line width=0.7pt,scale=0.75]\coordinate (P) at (0,0);\coordinate (Q) at (3.4,0);\coordinate (R) at (1.119,1.598);\draw (P) -- (Q) -- (R) -- cycle;\mbrwmarke{(R)}{(P)}{(Q)}
\node[font=\small,anchor=south east] at (0.52,0.82) {$m$};\node[font=\small,anchor=south west] at (2.30,0.82) {$n$};\node[font=\small,anchor=north] at (1.7,-0.08) {$k$};\end{tikzpicture}}{%
\vspace{6mm}\linie{50mm}}
\fusshilfe{\mbox{2:~$m^2 + n^2 = k^2$}}
\end{pfaufg}

% 3 Formel mit Wurzel (Kette 2 s8 v1); rechter Winkel oben rechts
\begin{pfaufg}{}{3.}{}
Kreuze die Formel an, mit der man $g$ berechnet.
\gzwei{\begin{tikzpicture}[line width=0.7pt,scale=0.75]\coordinate (A) at (0,1.8);\coordinate (B) at (3.2,1.8);\coordinate (C) at (3.2,0);\draw (A) -- (B) -- (C) -- cycle;\mbrwmarke{(B)}{(C)}{(A)}
\node[font=\small,anchor=south] at (1.6,1.88) {$e$};\node[font=\small,anchor=west] at (3.28,0.9) {$f$};\node[font=\small,anchor=north east] at (1.6,0.88) {$g$};\end{tikzpicture}}{%
\pfkreuzzeile{$g = \sqrt{e^2 + f^2}$\hspace{2.4em}$\square$\ $g = e^2 + f^2$}
\pfkreuzzeile{$g = \sqrt{e^2 - f^2}$\hspace{2.4em}$\square$\ $g = e + f$}}
\fusshilfe{\mbox{3:~$g = \sqrt{e^2 + f^2}$}}
\end{pfaufg}

% 4 kurzes Päckchen, krumme Zahlen: Kette 2 s1 v4, s4 v2, s5 v2, s3 v2
\begin{pfaufg}{}{4.}{}
Berechne die Hypotenuse $c$. Runde in d) auf eine Stelle nach dem Komma.\par\medskip
\gzwei{\begin{minipage}{66mm}
\pz{a)\ $36$ cm und $77$ cm}{=}{cm}
\pz{b)\ $4{,}5$ cm und $2{,}8$ cm}{=}{cm}
\pz{c)\ $2{,}4$ m und $70$ cm}{=}{cm}
\pz{d)\ $6$ m und $11$ m}{\approx}{m}
\end{minipage}}{\karo{3}}
\fusshilfe{\mbox{4:~a) 85 cm; b) 5,3 cm; c) 250 cm; d) $\sqrt{157} \approx 12{,}5$ m}}
\end{pfaufg}

% 5 Dreieck in anderer Lage (Kette 2 s2 v3): Hypotenuse oben waagerecht, ohne Hilfe
\begin{pfaufg}{}{5.}{}
Im Dreieck ist eine Seite mit ? markiert. Wie lang ist diese Seite?
\gzwei{\begin{tikzpicture}[line width=0.7pt,scale=0.7]\coordinate (A) at (0,2.48);\coordinate (B) at (5,2.48);\coordinate (C) at (2.16,0);\draw (A) -- (B) -- (C) -- cycle;\mbrwmarke{(C)}{(A)}{(B)}
\node[font=\small,anchor=south] at (2.5,2.56) {?};\node[font=\small,anchor=north east] at (1.0,1.2) {$48$ cm};\node[font=\small,anchor=north west] at (3.68,1.2) {$55$ cm};\end{tikzpicture}}{%
\karo{3}\par\smallskip\hfill ? $=$\linie{18mm}cm}
\fusshilfe{\mbox{5:~73 cm}}
\end{pfaufg}

% 6 Diagonale im Rechteck (Kette 2 s9 v2): ohne Skizze, das Dreieck findet der Schüler selbst
\begin{pfaufg}{}{6.}{}
Ein Rechteck ist $3{,}5$~m lang und $1{,}2$~m breit. Wie lang ist seine Diagonale?
\karo{3}\par\smallskip\hfill Diagonale:\linie{18mm}m
\fusshilfe{\mbox{6:~3,7 m}}
\end{pfaufg}

% 7 Teil eines Originals: Maße nur in der Skizze (2025 · 4 · a, Winkel weg)
\begin{pfaufg}{nach P10 ’25}{7.}{}
Familie Yücel baut vor der Haustür eine Rampe für Rollstühle. Berechne die Länge $x$ der Rampe.
\gzwei{\begin{tikzpicture}[line width=0.7pt,scale=0.62]\coordinate (P1) at (0,0);\coordinate (P2) at (5.95,0);\coordinate (P3) at (5.95,0.56);\draw (P1) -- (P2) -- (P3) -- cycle;\mbrwmarke{(P2)}{(P1)}{(P3)}
\node[font=\small,anchor=north] at (2.97,-0.08) {$170$ cm};\node[font=\small,anchor=west] at (6.03,0.28) {$16$ cm};\node[font=\small,anchor=south] at (2.9,0.32) {$x$};
\node[font=\scriptsize,mbgrau,anchor=north west] at (0,-0.5) {nicht maßstabsgerecht};\end{tikzpicture}}{%
\karo{3}\par\smallskip\hfill $x \approx$\linie{18mm}cm}
\fusshilfe{\mbox{7:~$x \approx 170{,}8$ cm}}
\end{pfaufg}

% 8 Überstand (Kette 2 s11 v2), ohne Hilfe; Draufsicht auf den Kistenboden
\begin{pfaufg}{}{8.}{}
Ein Stab liegt schräg in einer Kiste und ragt an einer Ecke hinaus. Wie lang ist der Stab?
\gzwei{\begin{tikzpicture}[line width=0.7pt,scale=0.72]
\draw (0,0) rectangle (2.4,2.52);
\draw[line width=1.8pt] (0,0) -- (2.617,2.748);
\mbrwmarke{(2.4,0)}{(0,0)}{(2.4,2.52)}
\node[font=\small,anchor=north] at (1.2,-0.08) {$40$ cm};
\node[font=\small,anchor=west] at (2.48,1.26) {$42$ cm};
\node[font=\small,anchor=south east] at (2.55,2.62) {$5$ cm};
\node[font=\scriptsize,mbgrau,anchor=north west] at (0,-0.5) {von oben gesehen};
\end{tikzpicture}}{%
\karo{4}\par\smallskip\hfill Stab:\linie{18mm}cm}
\fusshilfe{\mbox{8:~$58 + 5 = 63$ cm}}
\end{pfaufg}

% 9 Fehler finden (Kette 5 s1 v2): Wurzel aus jedem Summanden einzeln; „rechne richtig“ gestrichen (übt nur Nr. 4 noch einmal)
\begin{pfaufg}{}{9.}{}
Nora hat vier Hypotenusen berechnet. Genau eine Rechnung ist falsch. Kreuze sie an.\par\smallskip
{\small
\kk{}(1)\ Katheten $16$ cm und $30$ cm:\quad $c = \sqrt{256 + 900} = \sqrt{1\,156} = 34$ cm\par
\kk{}(2)\ Katheten $16$ cm und $63$ cm:\quad $c = \sqrt{256 + 3\,969} = \sqrt{4\,225} = 65$ cm\par
\kk{}(3)\ Katheten $21$ cm und $28$ cm:\quad $c = \sqrt{21^2 + 28^2} = 21 + 28 = 49$ cm\par
\kk{}(4)\ Katheten $28$ cm und $45$ cm:\quad $c = \sqrt{784 + 2\,025} = \sqrt{2\,809} = 53$ cm\par}
\fusshilfe{\mbox{9:~(3)}}
\end{pfaufg}

% 10 Prüfungshöhe (Kette 2 s15 v2, eigene Variante zu 2020 · 7 · a)
\begin{pfaufg}{}{10.}{}
Ein Boot fährt von Boje $A$ nach $C$ und dann rechtwinklig nach $B$. Berechne die Entfernung von $A$ nach $B$. Runde auf eine Stelle nach dem Komma.
\gzwei{\begin{tikzpicture}[line width=0.7pt]\coordinate (A) at (0,0);\coordinate (C) at (3.76,0);\coordinate (B) at (3.76,1.76);\draw (A) -- (C) -- (B);\draw[dashed] (A) -- (B);\mbrwmarke{(C)}{(A)}{(B)}
\foreach \p in {A,B,C}{\fill (\p) circle (1.6pt);}
\node[font=\small,anchor=east] at (A) {$A$};\node[font=\small,anchor=west] at (C) {$C$};\node[font=\small,anchor=west] at (B) {$B$};
\node[font=\small,anchor=north] at (1.88,-0.08) {$47$ m};\node[font=\small,anchor=west] at (3.84,0.88) {$22$ m};\end{tikzpicture}}{%
\karo{3}\par\smallskip\hfill $\overline{AB}^2 =$\linie{22mm}\quad $\overline{AB} \approx$\linie{18mm}m}
\fusshilfe{\mbox{10:~$\overline{AB}^2 = 2\,693$, $\overline{AB} \approx 51{,}9$ m}}
\end{pfaufg}

\blattende{2025 \pkt{} 4 \pkt{} a}{Vorher: Die lange Seite \pkt{} P10 \pkt{} 1\quad\textbullet\quad Weiter: Die lange Seite \pkt{} P10 \pkt{} 3}
\end{document}
